\documentclass[11pt,compress,t,notes=noshow, aspectratio=169, xcolor=table]{beamer}

\usepackage{../../style/lmu-lecture}
% Defines macros and environments
\input{../../style/common.tex}

\title{Applied Machine Learning}
% \author{LMU}
%\institute{\href{https://compstat-lmu.github.io/lecture_iml/}{compstat-lmu.github.io/lecture\_iml}}
\date{}

\begin{document}

\newcommand{\titlefigure}{figure/performance_vs_interpretability}
\newcommand{\learninggoals}{
\item What to watch out for in practice
\item Pros and cons of performance metrics
\item Advanced evaluation practices}

\lecturechapter{Applied Performance Evaluation}
\lecture{Applied Machine Learning}


\begin{frame}{General Things to Consider in Practice}

\begin{itemize}
\setlength\itemsep{2em}
    \item Stakeholders need to understand metrics. Stakeholders are (groups of) people that share an interest in the successful completion of the use case. 
    \item Is the interpretation important? Can the results be communicated to the stakeholders? Depending on the technical expertise of stakeholders or how much time can be invested in understanding the results, performance metrics shall be chosen and evaluated carefully.
    \item What matters for the specific usecase, e.g., outlier sensitivity, cost sensitivity?
    \item What are computational limitations (e.g., time constraints)? What computational resources are at our disposal?
    \item Are we developing a model that answers the right question? What a model optimizes does not necessarily overlap with what we are trying to achieve in a specific use case.
\end{itemize}
\end{frame}

\begin{frame}{Answering the Wrong Question}
% A. W. Kimball (1957) Errors of the Third Kind in Statistical Consulting, Journal of the American Statistical Association, 52:278, 133-142
\begin{itemize}
\setlength\itemsep{1em}
    \item An approximate answer to the right question is worth much more than the right answer to the wrong question.
    \item Pay attention to but do not get lost in technical details. You might end up developing a model that answers the wrong question or optimizes the wrong metric.
    \item \textbf{Example:} An online retail company aims to increase revenue by displaying advertisements with different products, hoping that the customer is going to purchase additional products. The company measures how often the user clicks / taps on those advertisements. A model is trained to predict click rates based on features such as the product category, price, age of the customer, etc. However, this model does not actually measure additional revenue from advertisements, as many customers will click on them by accident or not go through with the purchase.
\end{itemize}
    
\end{frame}
\begin{frame}{What Affects Performance?}
% https://vuquangnguyen2016.files.wordpress.com/2018/03/applied-predictive-modeling-max-kuhn-kjell-johnson_1518.pdf

There is a multitude of technical aspects we need to consider:
\vspace{1em}
\begin{itemize}
\setlength\itemsep{1em}
    \item \textbf{Model type:} Some models are or are not suited for certain tasks. For instance, tree-based models are bad at predicting linearly shaped data.
    \item \textbf{Non-informative features:} Including features that are not correlated with the target increases the risk of modeling noise.
    \item \textbf{Pre-processing of features:} Features can or have to be processed to improve performance, depending on the chosen model.
    \item \textbf{Measurement errors / noise in the outcome or features:} Garbage in = garbage out. Our performance depends on the quality of model inputs.
    \item \textbf{Sample size:} The more data we have at our disposal, the better our algorithm is able to learn from it.
    \item \textbf{What kind of data?:} It also matters what kind of data we have. Increasing the size of the data set is not sufficient. The data need to be diverse as well.
\end{itemize}
\end{frame}


\begin{frame}{Calibration versus Discrimination}

\begin{itemize}
\setlength\itemsep{2em}
    \item Calibration indicates the agreement between estimated probabilities and actual observed probabilities (through observed events).
    \item Discrimination is a model's ability to correctly separate observations into specific groups.
    \item Models can over-/underestimate probabilities (poor calibration) and still separate classes (good discrimination) or vice versa.
    \item Consider a model that predicts a patient's disease risk, where a patient is classified as a high-risk patient if disease risk exceeds 20\%. Patient A has a true risk of 5\% while patient B has a true risk of 30\%. The model predicts roughly the same risk with 19\% for A and 21\% for B. It is poorly calibrated but still separates both patients accurately into low-risk and high-risk groups (good discrimination).
    
\end{itemize}

\end{frame}

\begin{frame}{Selecting an Established Metric}

\begin{figure}
    \centering
    \includegraphics[width = 0.85\textwidth]{slides/04-perf-eval/figure/ontology_perf_metrics.png}
\end{figure}
\end{frame}

\begin{frame}{Selecting an Established Metric}

\begin{figure}
    \centering
    \includegraphics[width = 0.85\textwidth]{slides/04-perf-eval/figure/classification_of_metrics.png}
\end{figure}
\end{frame}


% \begin{frame}{Established Performance Metrics}

% \begin{itemize}
% \setlength\itemsep{2em}
%     \item List and motivate established performance metrics here
%     \item E.g., for \textbf{classification}: accuracy, true and false positive rate, precision and recall, F-score, AUC, Brier score
%     \item E.g., for \textbf{regression}: MSE, RMSE, MAE, R2
%     \item Describe typical use case and why metric works
%     \item Example: Brier Score rewards model confidence (i.e. give a high probability to a true outcome) and punishes the reverse.
% \end{itemize}
% \end{frame}

\begin{frame}{Advanced Performance Metrics}

\begin{itemize}
\setlength\itemsep{2em}
    \item Dive into esoteric stuff here: Youden, Cohens-Kappa, Matthews correlation
    \item Motivate each advanced metric, do not just list formula
    \item Example: "MCC produces a high score only if the prediction obtained good results in all of the confusion matrix categories (TP, FP, TN, FN), proportionally both to the size of positive elements and the size of negative elements in the dataset"
\end{itemize}
\end{frame}

% \begin{frame}{Lift Curves}

% \begin{itemize}
%     \item Lift charts (Ling and Li 1998) are a visualization tool for assessing the ability of a model to detect events in a data set with two classes. 
%     \item They rank the samples by their scores and determine the cumulative event rate as more samples are evaluated. In the optimal case, the M highest-ranked samples would contain all M events.
% \end{itemize}
% \end{frame}

\begin{frame}{Lift Charts}

\begin{itemize}
\setlength\itemsep{1em}
    \item Lift charts visualize a model's ability to detect events.
    \item Samples are ranked according to their scores (probability for an event to take place). We then evaluate the cumulative event rate.
    \item In the optimal case, we would suspect the $m$ highest ranked samples to contain all $m$ events.
    \item For a non-informative model, the fraction $\frac{k}{n}$ of the highest ranked samples contains $k$ events on average.
    \item For a more informative model, the fraction $\frac{k}{n}$ of the highest ranked samples contains more than $k$ events.
    \item The \textbf{lift} corresponds to the amount of events detected by a model above a completely random selection of samples.
\end{itemize}
\end{frame}



\begin{frame}{Lift Chart Example}
\begin{columns}
\begin{column}{0.5\textwidth}
\begin{itemize}
    \item We plot the cumulative event rate against the percentage of screened observations.
    \item Here, we face a 50\% event rate. For a perfect model, the 50\% highest ranked observations contain 50\% of all events.
    \item We create a lift curve for a non-informative model which runs along the diagonal line.
    \item Another model is able to perfectly detect events. After evaluating 50\% of the observations associated with the highest event probabilities, we have captured all events.
\end{itemize}
\end{column}
\begin{column}{0.5\textwidth}
\begin{figure}
    \centering
    \includegraphics[width = \textwidth]{slides/04-perf-eval/figure/lift_chart.png}
    \caption{Lift chart from Applied Predictive Modeling (Max Kuhn).}
\end{figure}
\end{column}

\end{columns}

\end{frame}


% \begin{frame}{Pseudo Code to Create Lift Chart}

% To construct a lift chart, we follow a general recipe:
% \begin{itemize}
% \setlength\itemsep{2em}
%     \item Predict target values for a sample that was not used during training with known target values.
%     \item Compute fraction of true events in the
% entire data set (= baseline event rate).
%     \item Rank data by scores (probability of event taking place).
%     \item For every unique event probability $\phi$, compute: 
%     $$\frac{\text{percentage of events in all observations with scores below $\phi$}}{\text{baseline event rate}}$$
% \end{itemize}
% \end{frame}

\begin{frame}{Calibration Curves}

\begin{columns}

\begin{column}{0.5\textwidth}
\begin{itemize}
\setlength\itemsep{2em}
    \item Calibration curves plot the true frequency of the positive label against its predicted probability for binned predictions.
    \item The horizontal axis represents the average predicted probability in each bin. The vertical axis indicates the the proportion of predicted positive labels.
    \item A well-calibrated classifier is associated with a diagonal clibration curve.
\end{itemize}
\end{column}
\begin{column}{0.5\textwidth}
\begin{figure}
    \centering
    \includegraphics{slides/04-perf-eval/figure/calibration_plot.png}
    \caption{Calibration plot from Applied Predictive Modeling (Max Kuhn).}
\end{figure}
\end{column}
\end{columns}

\end{frame}


\endlecture
\end{document}
